\section{Protocol}\label{sec:protocol}

\paragraph{Opponents.}
\begin{itemize}
\item \textbf{Global models} (\S\ref{sec:data}), each scored on four versions: raw and bias-corrected, under each of
the two run rules below. The bias correction is the trailing 30-day mean error at that station and hour of day,
taken from days before the valid time, with at least 10 such days required.
\item \textbf{The bias-corrected mean of IFS 0.25, AIFS, ICON and GFS}, issued before the forecast starts: a
classical model-output-statistics baseline.
\item \textbf{TimesFM 2.5} (200M parameters), \textbf{TimesFM 3.0} and \textbf{Chronos-2} (120M), zero-shot on each
station's last 512 hours. TimesFM 3.0 is used under its non-commercial research licence.
\item \textbf{Three simple baselines}: the same hour yesterday, the 14-day daily cycle, and the daily cycle plus the
latest anomaly.
\end{itemize}

\paragraph{Two rules for global models.}
\begin{itemize}
\item \textbf{STRICT} scores a global model on its latest run. That run is often issued after our forecast starts,
which favours the global model.
\item \textbf{ISSUED-BEFORE} scores it on its newest run published before our forecast starts: the 1-day-old run for
hours 0--16 and the 2-day-old run after. This is what a forecaster could actually have used.
\end{itemize}
Under either rule, \pwm \emph{beats} a global model on a variable only if its mean absolute error (MAE) is
significantly lower than both the raw and the bias-corrected forecast.

\paragraph{Significance.}
A difference counts only if the Diebold--Mariano test \citep{diebold1995dm} gives $p<0.05$ and a day-block bootstrap
\citep{kunsch1989bootstrap} 95\% interval excludes zero. Otherwise we call it a tie. Each comparison uses only the
hours where both forecasts and the observation exist.

\paragraph{Model choice on validation only.}
Three seeds of the default setting, P0, give a seed standard deviation $\sigma$. A candidate replaces P0 only if its
validation score is lower by more than $2\sigma$. The validation score is MAE in units of each variable's network
spread, averaged over the three variables. The best passing candidate is retrained with three fresh seeds for the
test. The same rule decides between hybrid designs. Appendix~\ref{app:variants} lists every candidate.
